\chapter{Análise Teórica}

Esta prática de laboratório é dividida em duas etapas principais. A primeira consiste na análise teórica do circuito para obtenção da função de transferência, magnitude e fase no domínio da frequência. A segunda etapa foca na montagem do circuito em protoboard e na coleta de dados experimentais utilizando o osciloscópio para diferentes frequências. Para a manipulação das equações algébricas e cálculos numéricos, utilizou-se software de computação simbólica.

\section{O Filtro de Realimentação Múltipla (MFB)}

O circuito analisado é um filtro passa-faixa de segunda ordem com configuração de realimentação múltipla (Multiple Feedback - MFB). Esta topologia utiliza um amplificador operacional para fornecer ganho e isolamento, enquanto os componentes passivos definem a seletividade da frequência. Considerando o modelo do amplificador operacional ideal e aplicando a Lei de Kirchhoff das Correntes (LKC) no domínio da frequência ($s$), a função de transferência é definida por:

\begin{equation}
	H(s) = \frac{V_o(s)}{V_i(s)} = \frac{-\frac{1}{R_1 C} s}{s^2 + \frac{2}{R_3 C} s + \frac{R_1 + R_2}{R_1 R_2 R_3 C^2}}.
	\label{eq:transfer_function}
\end{equation}

Para a análise em regime permanente senoidal, faz-se a substituição $s = j\omega$. A magnitude da resposta em frequência é dada por:

\begin{equation}
	|H(j\omega)| = \frac{\frac{\omega}{R_1 C}}{\sqrt{ \left( \frac{R_1 + R_2}{R_1 R_2 R_3 C^2} - \omega^2 \right)^2 + \left( \frac{2\omega}{R_3 C} \right)^2 }}.
	\label{eq:magnitude}
\end{equation}

A fase do sistema (em graus) é expressa como:

\begin{equation}
	\phi(\omega) = 180^\circ - 90^\circ - \arctan \left( \frac{\frac{2\omega}{R_3 C}}{\frac{R_1 + R_2}{R_1 R_2 R_3 C^2} - \omega^2} \right).
	\label{eq:phase}
\end{equation}

\section{Parâmetros e Valores Teóricos}

Para a validação do modelo, foram utilizados os seguintes valores nominais de componentes:
\begin{itemize}
	\item $C = 100$~nF
	\item $R_1 = 470$~$\Omega$
	\item $R_2 = 470$~$\Omega$
	\item $R_3 = 47$~k$\Omega$
\end{itemize}

A Tabela \ref{tab:valores_teoricos} apresenta a síntese dos valores calculados para as frequências limites e a frequência central de projeto.

\begin{table}[H]
	\centering
	\caption{Valores teóricos de magnitude e fase calculados.}
	\label{tab:valores_teoricos}
	\begin{tabular}{l|c|c}
		\hline
		\textbf{Frequência (Hz)} & \textbf{Magnitude $|H(j\omega)|$} & \textbf{Fase (Graus)} \\ \hline
		40 & 0,5947 & -90,68 \\ [3pt]
		480 (Pico/Ressonância) & 49,9733 & 178,13 \\ [3pt]
		11000 & 0,3084 & 90,35 \\ [3pt]
		\hline
	\end{tabular}
	\legend{Fonte: Elaborado pelos autores.}
\end{table}

Estes valores servem como referência para a comparação com os dados coletados experimentalmente, permitindo avaliar a precisão do modelo teórico em relação ao comportamento real do circuito montado em bancada.